\section{Conclusion}

\noindent
The Stefan-Boltzmann law has been validated for two different sources of thermal radiation, namly the Leslie's cube and a simulated blackbody.\\
The fit exponents of the different surfaces of the Leslie's cube turned out to be $b_{white} = 3.770 \pm 0.059$ for the white side,
$b_{black} = 4.043 \pm 0.081$ for the black side, $b_{polished} = 2.8 \pm 2.3$ for the polished nickel-plated side and $b_{matte} = 3.94 \pm 0.21$ for
the matte nickel-plated side. Exept the polished side the values for the fit exponents were plausible. The problem with the polished side was the missing
shielding from environmental radiation. The fit exponent of the blackbody resulted in $b_{blackbody} = 4.68 \pm 0.16$ with is also a bit off, but is still in
the same magnitude. These fit exponents suggest, that the Stefan-Boltzmann law is valid for both the Leslie's cube and the blackbody. \\

For different materials placed in the pathway between the thermopile and the source of thermal radiation, the transmittance can be calculated relative to
a free path. These resulted in $t_{glass} = (2.27 \pm 0.48)\%$ for window glass, $t_{silicon} = (46.82 \pm 0.48)\%$ for silicon and 
$t_{sodium chloride} = (73.32 \pm 0.48)\%$ which all match with the theoretical considerations. \\
According to Wien's law the maximal thermal radiation wavelength can be calculated. $\lambda_{max,0} = (9.885 \pm 0.033) \mu m$ for the room temperature
and $\lambda_{max,leslie} = (8.276 \pm 0.020 \mu m)$ for the Leslie's cube. \\
If the blackbody is not observed head-on, but at an angle, the intensity is expected to follow Lambert's cosine law. The thermal radiation emitted by the
oven was found to approximately follow this law especially for small angles.
